% Options for packages loaded elsewhere
\PassOptionsToPackage{unicode}{hyperref}
\PassOptionsToPackage{hyphens}{url}
%
\documentclass[
]{article}
\usepackage{amsmath,amssymb}
\usepackage{lmodern}
\usepackage{iftex}
\ifPDFTeX
  \usepackage[T1]{fontenc}
  \usepackage[utf8]{inputenc}
  \usepackage{textcomp} % provide euro and other symbols
\else % if luatex or xetex
  \usepackage{unicode-math}
  \defaultfontfeatures{Scale=MatchLowercase}
  \defaultfontfeatures[\rmfamily]{Ligatures=TeX,Scale=1}
\fi
% Use upquote if available, for straight quotes in verbatim environments
\IfFileExists{upquote.sty}{\usepackage{upquote}}{}
\IfFileExists{microtype.sty}{% use microtype if available
  \usepackage[]{microtype}
  \UseMicrotypeSet[protrusion]{basicmath} % disable protrusion for tt fonts
}{}
\makeatletter
\@ifundefined{KOMAClassName}{% if non-KOMA class
  \IfFileExists{parskip.sty}{%
    \usepackage{parskip}
  }{% else
    \setlength{\parindent}{0pt}
    \setlength{\parskip}{6pt plus 2pt minus 1pt}}
}{% if KOMA class
  \KOMAoptions{parskip=half}}
\makeatother
\usepackage{xcolor}
\setlength{\emergencystretch}{3em} % prevent overfull lines
\providecommand{\tightlist}{%
  \setlength{\itemsep}{0pt}\setlength{\parskip}{0pt}}
\setcounter{secnumdepth}{-\maxdimen} % remove section numbering
\ifLuaTeX
  \usepackage{selnolig}  % disable illegal ligatures
\fi
\IfFileExists{bookmark.sty}{\usepackage{bookmark}}{\usepackage{hyperref}}
\IfFileExists{xurl.sty}{\usepackage{xurl}}{} % add URL line breaks if available
\urlstyle{same} % disable monospaced font for URLs
\hypersetup{
  hidelinks,
  pdfcreator={LaTeX via pandoc}}

\author{Dietmar Gerald Schrausser}
\date{09.05.2026}


\begin{document}

\hypertarget{foliumblat-iiir}{%
\section{Foliu{[}m{]}/Blat IIIr}\label{foliumblat-iiir}}

A \emph{contemporary} English translation is presented from comparing
the existing Latin, Early New High German (ENHG) and English
(translation) texts for better comprehensibility. Especially since the
ENHG text is difficult to understand compared to modern High German, as
a \emph{profound} knowledge of German (as a mother tongue) as well as
already aknowledge of various \emph{dialects} is required.

\hypertarget{of-the-work-of-the-second-other-day}{%
\subsection{Of the Work of the Second (other)
Day}\label{of-the-work-of-the-second-other-day}}

\begin{quote}
\textbf{S}Ecundo die dixit deus. Fiat firmamentum in medio
aquaru{[}m{]}: et diuidat aquas ab aquis: vocauitq{[}ue{]} firmamentum
celum.
\end{quote}

Let us hear the prophet speak of celestial mysteries. God created a
\emph{wheel-shaped}\footnote{`rotundum' and `gescheibelt', disc-shaped.},
movable \emph{firmament}, in which other \emph{perceptible}\footnote{`sensibilia'
  and `empfintliche'.} things were located. From frozen water, it
solidified into a crystal and in this, the stars\footnote{`sidera' and
  `gestirne'.} were attached.\footnote{c.f. the `Speculum maius' (De
  Beauvais,
  \href{https://nbn-resolving.org/urn:nbn:de:hbz:061:1-20612}{1964},
  \emph{Speculum historiale}, Book I, Chapter 1), which describes the
  physical nature of the firmament and the \emph{waters above the
  heavens}, strongly influenced by Honorius Augustodunensis' `Imago
  Mundi'
  (\href{https://archive.org/detail/DHonoriiAugustudunensisLibriSeptem/page/n6/mode/1up}{1532}).}
The axes of the sphere of heaven, to which the stars are fixed, should
be \emph{two}. These, the (i) northern (\emph{Boreas}\footnote{Greek god
  of the north wind and personification of winter.}) and the (ii)
southern (\emph{Notus}\footnote{\emph{Notus}, the Greek god and
  personification of the south wind or \emph{nothus},
  \emph{illegitimate}, \emph{false}, \emph{mixed}, \emph{hybrid}.}),
\emph{roll} from east to west at such a speed that this would lead to
the collapse of the world were it not for the planets slowing this
movement in the opposite direction. Furthermore, the Creator of the
world tempered the nature of the heavens with waters, so that the
conflagration of the heavens would not set the lower elements
ablaze.\footnote{from the `Imago Mundi' (Augustudunensis,
  \href{https://archive.org/detail/DHonoriiAugustudunensisLibriSeptem/page/n6/mode/1up}{1532},
  \emph{De Coelo}, Book I, Chapter V), one of the most popular
  encyclopedias of the Middle Ages.}

{[}\ldots{]}\footnote{`Etlich lerer nit allein auß den unsern sunder
  auch auß den hebreyschen vn{[}d{]} andern haben geglawbt', Several
  teachers, not only ours but also among the Hebrews and others,
  believed that\ldots{}} \emph{Above}\footnote{`Supra', previously
  vs.~`uber'.} the (I) nine celestial spheres, that is, (i) the seven
planets, (ii) the eighth sphere, the \emph{unerring}\footnote{`inerrantem'
  and `die unirrende{[}n{]}'.}, (iii) the ninth orbit, comprehensible by
reason but not by the senses and the first among the \emph{moved}
bodies, it is assumed that (II) the tenth heaven is unchanging, at rest
and still, since it does not participate in any movement.

Not only our people held this view, in particular the younger
\textbf{`Strabo'} and \textbf{Bede}, but also many Hebrews and
furthermore some philosophers and mathematicians. It suffices to mention
two who attested to this: \textbf{Abraham}, the greatest Spanish
astrologer and \textbf{Isaac}, the philosopher. Here, \emph{Isaac}
interprets the tenth orbit as designated by \textbf{Ezekiel} through the
\emph{sapphire in the form of a throne}; the color of the sapphire
symbolizes the brilliance of light and the analogy\footnote{`similitudo'
  and `gleichnus'.} with a throne symbolizes immobility.\footnote{from
  the `Heptaplus' (Mirandola \& Mirandola,
  \href{https://doi.org/10.3931/e-rara-61217}{1601}, Book II, Chapter
  1), discusses the structure of the heavens and the existence of a
  tenth sphere (of the \emph{Empyreum} or the immobile \emph{first
  intelligence}) beyond the traditional nine spheres.}

⋅

But let us return to \emph{Moses}, who \textbf{\emph{separated the
waters from the waters in the midst of the firmament}}, resulting in a
threefold division of the \emph{sublunar} bodies. Several are located
(i) above the middle layer of air. In the uppermost part are its
elements and purest fire, entirely the \emph{ether}\footnote{`etheris'.},
clean, pure\footnote{`immixta' and `unvermischt'.} and true\footnote{`legitima'
  and `richtige'.} elements. Others are located (ii) below the middle
region\footnote{`meditulliu{[}m{]}' and `mitteln fürscheinenden stat'.}
of the air itself, that is, with us, where there is no pure element (and
also no pure perceptible element), but exclusively
\emph{mixtures}\footnote{`mixta' and `gemischt'.}, due to the
sediment-containing\footnote{`feculenta'.} and denser part of the
earthly body.\footnote{c.f. the `Calumniatorem Platonis' (Bessarion,
  \href{https://books.google.com/books?id=gPJCAAAAcAAJ}{1503}, Book IV,
  Chapter 6), a Neoplatonic interpretation of Genesis 1:6, which
  explains the \emph{division of the waters} by a threefold division of
  the sublunar bodies.}

The (iii) intermediate layer of air, \emph{also} called
\emph{firmament}, from which birds originate, being introduced by the
firmament of the heaven in flight, is the region in which the
\emph{etheric}\footnote{`sublimes'.} impressions such as rain, snow,
lightning, thunder, comets, or the like appear. \emph{This} firmament
(iii) therefore rightly separates the higher elements (i) from the lower
ones (ii), not only because of its position but also because of its
natural properties of different \emph{waters}. Above\footnote{`Sup{[}ra{]}'
  and `darob'.} it are the pure elements, below them, \emph{completely}
mixed in elemental simplicity\footnote{`simplicitate' and
  `slechtigkeit'.} separated\footnote{`discedunt' and `obgesünderet'.}.
Thus, he named this firmament, derived from koylon (concave), as the
\emph{heaven}\footnote{`celum' (koylon) and `himel'.}, since it
encompasses all perceptible\footnote{`sensibilia' and `empfintliche'.}
and invisible things.\footnote{c.f. the `Epitoma in Almagestum
  Ptolemaei' (Regiomontanus,
  \href{https://doi.org/10.3931/e-rara-528}{1496}, \emph{Praefatio}),
  reflects a medieval synthesis of biblical cosmology and Aristotelian
  physics.}

\hypertarget{references}{%
\subsection{References}\label{references}}

Augustudunensis, H. (1532). \emph{D. Honorii Augustudunensis Presbyteri
Libri Septem}. Basileae.
\url{https://archive.org/detail/DHonoriiAugustudunensisLibriSeptem/page/n6/mode/1up}

Bessarion. (1503). \emph{In Calumniatorem Platonis \ldots: Libri IV}.
Venetiis: Aldus Manutius.
\url{https://books.google.com/books?id=gPJCAAAAcAAJ}

De Beauvais, V. (1964). \emph{Speculum maius}. Graz: Akad. Druck- u.
Verl.-Anst. \url{https://nbn-resolving.org/urn:nbn:de:hbz:061:1-20612}

Mirandola, G., \& Mirandola, G. F. (1601). \emph{Ioannis Pici,
Mirandulae \ldots{} opera quae extant omnia \ldots{}} Basileae: per
Sebastianum Henricpetri. \url{https://doi.org/10.3931/e-rara-61217}

Regiomontanus. (1496). \emph{Epytoma Ioannis de Monte Regio in
almagestum Ptolemei}. Venetiis: Johannes Hamman.
\url{https://doi.org/10.3931/e-rara-528}

\end{document}
